\documentclass[11pt,a4paper]{article}
\usepackage[T1]{fontenc}\usepackage[utf8]{inputenc}
\usepackage{lmodern,microtype,amsmath,amssymb,amsthm,mathtools}
\usepackage[margin=25mm]{geometry}
\usepackage{booktabs,array}\usepackage[hyphens]{url}\usepackage[hidelinks]{hyperref}
\setlength{\emergencystretch}{3em}\allowdisplaybreaks
\newtheorem{theorem}{Theorem}[section]
\newtheorem{lemma}[theorem]{Lemma}
\newtheorem{proposition}[theorem]{Proposition}
\newcommand{\ph}{\Phi}\newcommand{\ord}{\operatorname{ord}}
\hypersetup{pdftitle={Karmakar--Jana Conjecture 2: a root-of-unity proof candidate},pdfauthor={Reiwa}}
\title{Karmakar--Jana Conjecture 2\\[4pt]\large A global root-of-unity proof candidate, with a failed new-method gate}
\author{Reiwa: tournament round 11}
\date{16 November 2026}
\begin{document}\maketitle
\begin{abstract}
We examine Conjecture 2 of Karmakar and Jana, arXiv:2610.11778v1,
which strengthens their equation (1.18) from $\Phi_m(q)^5$ to
$[m]\Phi_m(q)^4$ for every odd $m>1$ and every permitted shift.
This report supplies a full proof candidate based on their proved Theorem 2,
a complete-block reduction, and explicit cancellation of the poles in the
specialized harmonic correction. Composite divisors and all endpoints are
included. Independent verification is pending. The block method has clear
antecedents, including the same authors' arXiv:2610.15050v1 and our round 7;
the new-method gate G5 therefore fails. No priority or new-method claim is made.
\end{abstract}

\section{Exact original assertion and status}
The target is \emph{Conjecture 2}, Section 4, page 11 of \cite{KJ}.
Its reference is equation (1.18), page 5, not (1.14) and not the
base-$q^4$ conjecture treated in \cite{G2}. Both source pages were visually
checked. The official arXiv record inspected on 16 November 2026 lists
only version 1, submitted 8 October 2026. A bounded search found no later
resolution of this exact conjecture; this is not a certification of open status.

Let $q$ be indeterminate, $m>1$ odd, $h=(m-1)/2$, and
$s\in\{0,1,\ldots,h\}$. Set
\[
 [a]=\frac{1-q^a}{1-q},\qquad
 (z;q^2)_k=\prod_{i=0}^{k-1}(1-zq^{2i}),\qquad
 R_k=\frac{(q;q^2)_k}{(q^2;q^2)_k}\quad(k\geq0).
\]
Empty products are one and empty sums are zero. Define
\begin{align}
 L_{m,s}&=\sum_{k=s}^{h}[4k+1]R_{k-s}R_{k+s}R_k^4q^k,\label{eq:L}\\
 B_{m,s}&=[m+2s]R_s^4\frac{R_{h+s}}{R_{h-s}}q^{-h+2s},\label{eq:B}\\
 U_{s,j}&=R_j^2
 \left(\frac{(q^{2s+1};q^2)_j}{(q^{2s+2};q^2)_j}\right)^2q^{2j},\label{eq:U}\\
 H_{s,j}&=\sum_{i=1}^j\left\{
 \frac{q^{2(s+i)-1}}{[2(s+i)-1]^2}
 +\frac{q^{2(s+i)}}{[2(s+i)]^2}\right\},\label{eq:H}\\
 V_{m,s}&=B_{m,s}\sum_{j=0}^{h-s}U_{s,j}
 \left\{1-(2-q^m)[m]^2H_{s,j}\right\}.\label{eq:V}
\end{align}
These are exactly the two sides of (1.18), with its summation variable
renamed and its products compressed into $R_k$. In particular the fourth
powers in \eqref{eq:L} and \eqref{eq:B} and the two squared shifted factors
in \eqref{eq:U} have not been replaced by factorials with index $4k$ or $2j$.
The source's already-proved Theorem 2 states
\begin{equation}
 L_{m,s}\equiv V_{m,s}\pmod{\ph_m(q)^5}.\label{eq:known}
\end{equation}
The \emph{entire} conjecture is
\begin{equation}
 L_{m,s}\equiv V_{m,s}\pmod{[m]\ph_m(q)^4}
 \quad(m>1\text{ odd},\ 0\leq s\leq h).\label{eq:target}
\end{equation}
Congruences are in $\mathbb Q(q)$ after cancellation: the reduced denominator
must be coprime to the stated polynomial modulus. Establishing this
regularity is part of the argument, not an assumption about the displayed
uncancelled factors.

\section{Cyclotomic orders and the specialized right side}
Fix any divisor $\ell>1$ of $m$, not necessarily prime. Write
$t=(\ell-1)/2$ and let $\nu_\ell$ be the exponent of $\ph_\ell(q)$ in
a nonzero rational function. Every zero of $1-q^a$ for positive integer
$a$ is simple, and $\ph_\ell\mid1-q^a$ exactly when $\ell\mid a$.
For $k\geq0$, counting odd numerator and even denominator exponents gives
\begin{equation}
 \nu_\ell(R_k)=\left\lfloor\frac{k+t}{\ell}\right\rfloor
 -\left\lfloor\frac{k}{\ell}\right\rfloor
 =:\delta(k)
 =\begin{cases}0,&0\leq k\bmod\ell\leq t,\\1,&t<k\bmod\ell<\ell.
 \end{cases}\label{eq:delta}
\end{equation}
In particular every $R_k$ is regular at primitive $\ell$th roots, even
when its displayed denominator contains zero factors. Thus every left
summand in \eqref{eq:L} is regular there.

\subsection{Group first, then specialize}
One cannot substitute a root into $B_{m,s}$ and $U_{s,j}$ separately
without checking poles. The identity
\[
 \frac{(q^{2s+1};q^2)_j}{(q^{2s+2};q^2)_j}=\frac{R_{s+j}}{R_s}
\]
holds in $\mathbb Q(q)$ and yields the useful grouping
\begin{align}
 D_{m,s}&=[m+2s]R_s^2\frac{R_{h+s}}{R_{h-s}},\label{eq:D}\\
 W_{m,s,j}:=B_{m,s}U_{s,j}
 &=D_{m,s}R_j^2R_{s+j}^2q^{-h+2s+2j}.\label{eq:W}
\end{align}
Since $m/\ell$ is odd, write $h=A\ell+t$ with
$A=(m/\ell-1)/2$. For $u=s\bmod\ell$, equation \eqref{eq:delta} gives
\[
 \nu_\ell(D_{m,s})=
 \boldsymbol1_{u=0}+2\delta(s)+\delta(h+s)-\delta(h-s)=1.
\]
For clarity all cases are displayed:
\[
\begin{array}{c|rrrr|r}
 u&\boldsymbol1_{u=0}&2\delta(s)&\delta(h+s)&-\delta(h-s)&\text{sum}\\\hline
 0&1&0&0&0&1\\
 1\leq u\leq t&0&0&1&0&1\\
 t<u<\ell&0&2&0&-1&1
\end{array}
\]
Here $\ell\mid m+2s$ is equivalent to $u=0$, since $\ell$ is odd.
The negative power of $q$ in \eqref{eq:W} is a unit at all roots in question.
Consequently
\begin{equation}
 \nu_\ell(W_{m,s,j})=1+2\delta(j)+2\delta(s+j)\geq1.
 \label{eq:Worder}
\end{equation}

\subsection{The harmonic correction can have a nonzero limit}
Every summand in $H_{s,j}$ has a pole of order at most two, because
$1-q$ is nonzero at primitive $\ell$th roots. Since
$\nu_\ell([m]^2)=2$, the product $[m]^2H_{s,j}$ is regular. It follows
from \eqref{eq:Worder} that each \emph{grouped} right summand
\[
 W_{m,s,j}\{1-(2-q^m)[m]^2H_{s,j}\}
\]
is regular and has order at least one. Hence
\begin{equation}
 V_{m,s}(\zeta)=0
 \quad\text{for every primitive $\ell$th root $\zeta$ and every $\ell\mid m$, $\ell>1$}.
 \label{eq:Vzero}
\end{equation}
No interchange with an infinite sum or parameter limit is involved.

The cancellation above must not be replaced by the false assertion
$[m]^2H_{s,j}(\zeta)=0$. In fact, termwise evaluation of this finite sum gives
\begin{equation}
 \lim_{q\to\zeta}[m]^2H_{s,j}(q)
 =\sum_{\substack{2s+1\leq r\leq2(s+j)\\\ell\mid r}}
 \left(\frac{m}{r}\right)^2.\label{eq:Hlimit}
\end{equation}
For a divisible $r$, use $(1-q^m)/(1-q^r)\to m/r$ and $\zeta^r=1$;
all other terms tend to zero. For instance $m=9$, $s=0$, $j=2$, $\ell=3$
gives the limit $9$, so that the correction bracket tends to $-8$, not $1$.
It is the zero of $W$, not a zero of the bracket, that proves \eqref{eq:Vzero}.

\section{Complete blocks and the shifted left side}
\begin{lemma}[Exact root limits]
Let $\ell>1$ be odd, $t=(\ell-1)/2$, and $\zeta$ a primitive $\ell$th
root. For integers $a\geq0$ and $0\leq r<\ell$, put
$\beta_a=(1/2)_a/a!$. Then
\begin{equation}
 R_{a\ell+r}(\zeta)=
 \begin{cases}
 \beta_a R_r(\zeta),&0\leq r\leq t,\\
 0,&t<r<\ell.
 \end{cases}\label{eq:blockR}
\end{equation}
Every value means the regular rational-function limit after cancellation.
\end{lemma}
\begin{proof}
Separate the index product into blocks of length $\ell$. In block number
$b\geq0$, the unique numerator zero is $1-q^{(2b+1)\ell}$ and the unique
denominator zero is $1-q^{(2b+2)\ell}$. Their ratio tends to
$(2b+1)/(2b+2)$. After these factors are removed, the remaining numerator
and denominator exponents each run through all nonzero residues modulo
$\ell$: multiplication by $2$ permutes residues because $\ell$ is odd.
Their products therefore cancel at $\zeta$. The $a$ full blocks contribute
$\prod_{b=0}^{a-1}(2b+1)/(2b+2)=\beta_a$.
The remaining partial block has no denominator zero. If $r\leq t$ it also
has no numerator zero and equals $R_r(\zeta)$; otherwise it contains
one numerator zero. This proves \eqref{eq:blockR}, without assuming that
$\ell$ is prime or that nonzero residues are relatively prime to $\ell$.
\end{proof}

\begin{lemma}[Signed-shift reduction]
Under the notation above, let $h=A\ell+t$ and $0\leq s\leq h$.
Choose the nearest multiple of $\ell$ to $s$, writing
\begin{equation}
 s=b\ell+\varepsilon v,\qquad
 0\leq b\leq A,\quad 0\leq v\leq t,\quad
 \varepsilon\in\{-1,1\}.\label{eq:signed}
\end{equation}
For $v=0$ either sign works. Then
\begin{equation}
 L_{m,s}(\zeta)=
 \left(\sum_{a=b}^{A}\beta_{a-b}\beta_{a+b}\beta_a^4\right)
 L_{\ell,v}(\zeta).\label{eq:blockL}
\end{equation}
\end{lemma}
\begin{proof}
Write $k=a\ell+r$ with $0\leq r<\ell$. By \eqref{eq:delta}, a summand
can have a nonzero root value only if
$\delta(k)=\delta(k-s)=\delta(k+s)=0$. For $r\leq t$ the last two
conditions are precisely
\begin{equation}
 v\leq r\leq t-v.\label{eq:support}
\end{equation}
Indeed $r-v<0$ would have residue $\ell+r-v>t$, while $r+v>t$
would have residue in $(t,\ell)$; neither can survive. Conversely the
interval \eqref{eq:support} keeps both $r-v$ and $r+v$ in $[0,t]$.
When $2v>t$ it is empty, and both sides of \eqref{eq:blockL} vanish.

Otherwise the original bounds $s\leq k\leq h$ select exactly
$b\leq a\leq A$ with all $r$ in \eqref{eq:support}. To check the lower
endpoint also when $\varepsilon=-1$, any $a<b$ gives
$k\leq(b-1)\ell+t-v<b\ell-v=s$, whereas $a=b$ is allowed.
For $\varepsilon=1$, $a=b$ is allowed exactly because $r\geq v$.
The upper endpoint is allowed since $r\leq t-v\leq t$.
In \eqref{eq:signed}, $b\leq A$ follows from $s\leq A\ell+t$ and
$t<\ell/2$; if $\varepsilon=-1$ and $v>0$, necessarily $b\geq1$.
Thus no negative factorial index is introduced.

Within a surviving block,
\[
 k-s=(a-b)\ell+(r-\varepsilon v),\qquad
 k+s=(a+b)\ell+(r+\varepsilon v).
\]
There is no carry or borrow, and \eqref{eq:blockR} transforms their
product into $\beta_{a-b}\beta_{a+b}R_{r-v}(\zeta)R_{r+v}(\zeta)$.
Also $R_k^4$ contributes $\beta_a^4R_r(\zeta)^4$,
$\zeta^k=\zeta^r$, and $[4k+1](\zeta)=[4r+1](\zeta)$.
Consequently every surviving block is its displayed beta factor times
\[
 \sum_{r=v}^{t-v}[4r+1](\zeta)
 R_{r-v}(\zeta)R_{r+v}(\zeta)R_r(\zeta)^4\zeta^r.
\]
This equals $L_{\ell,v}(\zeta)$: its omitted indices
$t-v<r\leq t$ have $R_{r+v}(\zeta)=0$. Summing the blocks proves
\eqref{eq:blockL}.
\end{proof}

\section{Completion of the proposed global argument}
\begin{theorem}[Proof candidate for the complete original conjecture]
For all odd integers $m>1$ and all $0\leq s\leq(m-1)/2$,
equation \eqref{eq:target} holds in the reduced-denominator sense.
\end{theorem}
\begin{proof}
For any divisor $\ell>1$ of $m$, the preceding arguments establish
regularity of both sides at every primitive $\ell$th root.
Apply the \emph{proved} source theorem \eqref{eq:known} with its own
modulus parameter equal to $\ell$ and its shift equal to $v$:
\[
 L_{\ell,v}\equiv V_{\ell,v}\pmod{\ph_\ell^5}.
\]
Equation \eqref{eq:Vzero}, already proved directly and also applicable
when $m=\ell$, gives $V_{\ell,v}(\zeta)=0$. Thus
$L_{\ell,v}(\zeta)=0$, and \eqref{eq:blockL} gives
$L_{m,s}(\zeta)=0$. Together with \eqref{eq:Vzero}, this yields
\[
 L_{m,s}\equiv V_{m,s}\equiv0\pmod{\ph_\ell}
 \qquad(1<\ell\mid m).
\]
The argument uses Theorem 2 only in its established $\ph_\ell^5$ form,
not the conjectured stronger form. There is no induction on an unproved
assertion and no reliance on our round-7 Conjecture-1 candidate.

In characteristic zero the distinct cyclotomic polynomials are pairwise
coprime, and
\[
 [m]=\prod_{\substack{\ell\mid m\\\ell>1}}\ph_\ell(q).
\]
The regularity just proved means that the reduced denominators are
coprime to this product. We have therefore proved the first-order
congruence modulo $[m]$. Independently, \eqref{eq:known} supplies
order at least five at $\ph_m$. Taking the least common multiple,
\[
 \operatorname{lcm}([m],\ph_m^5)=[m]\ph_m^4,
\]
proves precisely \eqref{eq:target}. Neither $q=1$ nor a prime-only
specialization is being substituted for the original statement.
\end{proof}

The endpoint $s=h$ can also be checked without any congruence: the right
sum has only $j=0$, and both sides equal
$[4h+1]R_{2h}R_h^4q^h$. The endpoint $s=0$, the smallest modulus $m=3$,
and composite $m$ all fall within the proofs. For prime $m$, the target
modulus is already $\ph_m^5$; the substantive strengthening concerns
proper cyclotomic factors for composite $m$.

\section{Reproducible checks and deliberate negative controls}
The proof above is symbolic. The computations below are finite checks
for transcription, endpoints, cancellations, and multiplicity; they do
not establish the universal theorem.

\paragraph{Exact rational polynomial arithmetic.}
The supplied \texttt{checks.py} explicitly strips powers of $\ph_\ell$
from each $1-q^a$ before inverting units in
$\mathbb Q[q]/(\ph_\ell^e)$. It never inverts an uncancelled zero factor.
There are 232 successful cases in \texttt{exact\_results.json}:
20 fifth-order cases, comprising all shifts for $m=3,5,7,9,11$, and
212 first-order proper-divisor cases from
$m\in\{9,15,21,25,27,35,45\}$.
Each proper-divisor case verifies that the left and right sides vanish
separately. Replacing $2-q^m$ by $3-q^m$ gives a nonzero fifth-order
residual for $m=3,5,7$ at $s=0$, as required by three negative controls.

\paragraph{Independent finite-field jets.}
The supplied \texttt{jet\_check.py} builds the \emph{ungrouped original}
right factors of \eqref{eq:B}--\eqref{eq:U}, rather than importing
\eqref{eq:W} from the first program. It works with $q=\zeta+\epsilon$
over a prime field with an element of exact order $\ell$, factors the
individual simple zeros, and keeps five coefficients when $\ell=m$.
The primes exceed the largest exponents used, so no zero acquires an
accidental characteristic-$p$ multiplicity. The 396 successful cases
comprise 150 fifth-order jets and 246 first-order proper-divisor values,
for all shifts and relevant divisors of
\[
 m\in\{3,5,7,9,11,13,15,17,19,21,25,27,33,35,45\}.
\]
Four negative controls, $m=3,5,7,9$ and $s=0$, again detect the wrong
factor $3-q^m$. One root per finite field is tested, not every algebraic
root; the proof, rather than this test, covers them all.

The same independent program exhaustively checks the valuation table
and exact signed-shift support for all odd $3\leq m\leq301$, every
$\ell>1$ dividing $m$, and every allowed $s$: 33,213 cases.
The data, field primes, roots, residuals of negative controls, and
runtimes are recorded in \texttt{jet\_results.json}. The explicit limit
\eqref{eq:Hlimit} explains why first-order root checks alone could not
validate the correction coefficient.

\section{Prior art, remaining verification, and gate verdict}
The official records of \cite{KJ,G2}, the author's publication page
\cite{Jana}, and searches by exact identifier, title, and conjecture number
were inspected on 16 November 2026. The later paper \cite{G2} proves
Tang's Conjecture 4.1 and Wang--Xu's Conjecture 2 in a base-$q^4$ setting;
its Theorem 1.6 concerns $m\equiv1\pmod4$. It is not a stated resolution
of the present base-$q^2$ conjecture for all odd $m$. Nevertheless its
Lemma 2.1 and subsequent shifted root-block arguments are direct method
antecedents. The $q^3$ Andrews-transformation paper arXiv:2610.33249 was
also screened and is a different problem. A secondary aggregator that
conflates these source problems was not used to define the target.

The 190-source corpus snapshot at commit
\texttt{5887cdea4898f0a3aa5bf7a8e0829697bd2b0f4b} was searched for
\texttt{2610.11778}, \texttt{Karmakar}, the title fragment
\texttt{Third and Fifth Powers}, and \texttt{2610.15050}; no hits were
found. The reproducible literal audit is \texttt{audit.py} and
\texttt{corpus\_audit.json}. This limited search cannot certify semantic
nonduplication. In addition, our own round 7 used the same root-block
mechanism for the distinct Conjecture 1 in \cite{KJ}. The present direct
pole treatment avoids unjustified parameter specialization, but does not
make the underlying method substantively new.

\begin{center}\small
\begin{tabular}{p{.09\linewidth}p{.82\linewidth}}
\toprule
Gate&Verdict and scope\\\midrule
G1&Pass: an explicitly numbered, pre-existing Conjecture 2.\\
G2&Pass: all of (1.18), all odd $m>1$, all original shifts, and the original stronger modulus.\\
G3&Bounded evidence only: original still labels it open; no later exact resolution located; global open status unverified.\\
G4&Full proof candidate supplied, using an established theorem rather than a conjecture. Independent verification pending.\\
G5&Fail: known block strategy, also reused from round 7; no substantively new-method claim.\\
G6&Not certified: explicit source/corpus screening, but method overlap is known and exact priority remains unverified.\\\bottomrule
\end{tabular}
\end{center}

No mathematical step is intentionally left as a conjecture in the
candidate. The outstanding checks are independent scrutiny of the exact
source transcription, pole/grouping argument, signed-shift block bounds,
and final multiplicity combination, plus a broader historical search.
Even if all mathematical steps are accepted, this submission does not
pass all six tournament gates because G5 fails.

\begin{thebibliography}{9}
\bibitem{KJ} L.~Karmakar and A.~Jana,
\emph{$q$-Supercongruences Modulo the Third and Fifth Powers of a Cyclotomic Polynomial},
arXiv:2610.11778v1 (8 October 2026), Theorem 2, equation (1.18), p.~5;
Conjecture 2, Section 4, p.~11.
\url{https://arxiv.org/abs/2610.11778}.
\bibitem{G2} L.~Karmakar and A.~Jana,
\emph{Proof of Two Conjectural $q$-Supercongruences Related to Van Hamme's (G.2) Supercongruence},
arXiv:2610.15050v1 (10 October 2026), Theorems 1.5--1.6 and Section 2.
\url{https://arxiv.org/abs/2610.15050}.
\bibitem{Jana} A.~Jana, author publication list, inspected 16 November 2026.
\url{https://sites.google.com/view/arijitjana/publications}.
\bibitem{Corpus} QEEDDD,
\emph{Verifiable Mathematical Reasoning Corpus}, inspected snapshot at
commit \texttt{5887cdea4898f0a3aa5bf7a8e0829697bd2b0f4b}.
\url{https://github.com/QEEDDD/Verifiable-Mathematical-Reasoning-Corpus}.
\end{thebibliography}
\end{document}
